\ReviewerComment{Why was MPC-4 selected as the primary comparison baseline? In addition, the inference time of the DQN models and the decision time of the MPC method should be evaluated and compared on the same hardware platform under the same experimental conditions.}
\AuthorResponse{We now explain that four steps provide nontrivial look-ahead while retaining transparent exhaustive enumeration at every decision. We do not treat $H=4$ as globally optimal and retain the $H\in\{1,2,4,6\}$ planning-depth trade-off. We added an eight-step greedy rollout and a matched CPU benchmark for all six DQN objectives, MPC-4, and greedy rollout. Measurements use one CPU thread, batch size one, shared saved states, warm-up, and three timed repetitions, and report mean, median, and P95 decision time. The matched comparison appears on pp.~20--21 (Table~12).}
\RevisedExcerpt{Mean DQN decision time was 0.050--0.051~ms, compared with 2.538~ms for greedy rollout and 4.772~ms for MPC-4.}
\ManuscriptLocation{pp.~12--13, Section~3.4.4; p.~18, Table~10; pp.~20--21, Section~4.6 and Table~12.}

\ReviewerComment{Both the training and evaluation data used in this study are generated by a synthetic simulator, and the reported findings have not yet been validated using real satellite mission data. Although the authors describe the implemented parameter ranges and provide general engineering references for the synthetic simulator, the relationship between its parameter settings and actual satellite operating conditions remains insufficiently explained. The authors are encouraged to clarify how these parameter ranges correspond to representative satellite platforms or mission scenarios and to discuss how potential discrepancies may affect the reported findings. If feasible, validation using real mission data would further strengthen the practical relevance and generalizability of the study.}
\AuthorResponse{We expanded the parameter-provenance explanation and distinguish three categories. Link rates and computing power are engineering-scale anchors supported by cited public sources. Data-reduction ratios and task correlations are synthetic assumptions. Heat, energy, queue, utilization, bandwidth, and contact coordinates are normalized accounting states rather than a calibrated spacecraft digital twin. The Discussion now explains that mission telemetry and hardware-in-the-loop validation are required before operational transfer.}
\RevisedExcerpt{The parameters serve three distinct roles. Compute power and link rates are engineering-scale anchors supported by public NASA ranges, including SmallSat computing-system power and ground-network examples from 2.2 Mbit s$^{-1}$ S-band to 220 Mbit s$^{-1}$ X-band returns~\cite{nasa2026avionics,nasa2026ground}. Data-reduction factors and task correlations are synthetic workload assumptions; the reduction factors represent the bandwidth-reduction role of onboard image compression rather than an instrument-specific codec~\cite{ccsds122}. Heat, energy, queue, utilization, bandwidth, and contact coordinates are normalized accounting states rather than calibrated hardware variables. Table~\ref{tab:parameters} therefore defines a transparent engineering envelope, not a fitted flight distribution or spacecraft digital twin; mission-specific discrepancies in these scales and joint distributions may change absolute costs and potentially the method ordering.}
\ManuscriptLocation{p.~6, Section~3.1.2; p.~7, Table~3; p.~23, Sections~5.1--5.2; p.~24, Conclusions.}
